\subsubsection{Obtain the normal from Miller indices}

A crystal plane's normal is determined by reciprocal vectors. Using the convention in Chapter~\ref{ch:theory},
\begin{equation}
  \mathbf G_{hkl}=h\mathbf b_1+k\mathbf b_2+l\mathbf b_3,
  \qquad
  \hat{\mathbf n}=\frac{\mathbf G_{hkl}}{\|\mathbf G_{hkl}\|}.
  \label{eq:surface-miller-normal}
\end{equation}
The reciprocal vectors $\mathbf b_1,\mathbf b_2,\mathbf b_3$ have units $\text{\AA}^{-1}$; $h,k,l$ are dimensionless integers and $\hat{\mathbf n}$ is a unit normal. In a cubic cell, the (111) normal is parallel to $\mathbf a+\mathbf b+\mathbf c$. Here the lattice axes align with global axes, so it is also parallel to global $(1,1,1)$. A skewed cell generally lacks this correspondence: do not use Miller indices directly as a Cartesian adsorbate-orientation vector.

For a constructed slab with in-plane vectors $\mathbf a,\mathbf b$, check its actual normal and cell height using
\begin{equation}
  \hat{\mathbf n}_{\mathrm{slab}}=
  \frac{\mathbf a\times\mathbf b}{\|\mathbf a\times\mathbf b\|},
  \qquad H=|\mathbf c\cdot\hat{\mathbf n}_{\mathrm{slab}}|
  \label{eq:surface-normal-height}
\end{equation}
The height $H$ is the repeat distance between periodic slabs along the normal, and need not equal $\|\mathbf c\|$. First choose periodic images that form one continuous slab, then project the atomic positions onto the normal to measure the distance between the outermost layers. Record this distance separately from the requested minimum thickness and the thickness of the complete structural repeats used for construction.

\begin{table}[H]
\centering
\caption{Lengths used in surface construction.}
\begin{tabularx}{\linewidth}{lX}
\toprule
Quantity & Interpretation and checks\\
\midrule
Minimum slab thickness & A criterion for selecting complete repeat units, not a direct layer count.\\
Atomic-layer span & Normal distance between outermost layers; select one continuous slab first.\\
Cell height $H$ & Measured along the surface normal, Eq.~\eqref{eq:surface-normal-height}; includes the slab and periodic gap.\\
Added vacuum $\Delta$ & Length added along the third lattice vector by \ui{Add Vacuum}, not the final total vacuum.\\
\bottomrule
\end{tabularx}
\end{table}

For nonorthogonal input, also check the output normal, layer count, and periodic separation. Orientation handling can affect thickness conversion and candidate reduction; the minimum input thickness alone does not establish the result. The platinum example starts from a cubic cell and checks the generated slab.
